\documentclass[10pt]{article}
\usepackage[diary]{aditya}

\title{Diary entry for Brown-Representability}
\author{Aditya Pahuja}
\date{\today}

\begin{document}
\maketitle
\tableofcontents
\pagebreak

\section{$\Omega$-spectra represent cohomology theories}
\begin{definition}[$\Omega$-spectrum]
    \label{def:omega-spectrum}
    An $\Omega$-spectrum is a sequence of spaces
    $K_1$, $K_2$, \dots equipped with (weak) homotopy equivalences
    $\xi_n\colon K_n\to\Omega K_{n+1}$ for each $n$.
\end{definition}

\begin{example}
    [Eilenberg-MacLane spectra]
    Let $K_n = K(G,n)$ be the $n$th Eilenberg-MacLane space.
    For a given $n$ and any $k$,
    \begin{equation*}
	\pi_k(K_n)\cong \pi_{k+1}(K_{n+1})
	= \pi_k(\Omega K_{n+1}),
    \end{equation*}
    so $K_n\simeq\Omega K_{n+1}$
    because Eilenberg-MacLane spaces are unique up to homotopy.
    Choosing $\xi_n$ to be a homotopy equivalence
    from $K_n$ to $K_{n+1}$ thus gives an $\Omega$-spectrum
    for each abelian group $G$, often denoted $HG$
    and called the \vocab{Eilenberg-MacLane spectrum of $G$}.
\end{example}

\begin{theorem}[$\Omega$-spectra represent cohomology theories]
    \label{thm:omega-spectra-cohomology}
    Let $\{(K_n,\xi_n)\}$ be an $\Omega$-spectrum.
    The sequence of contravariant functors $H^n = [-,K_n]$
    from the category of pointed CW complexes
    to the category of abelian groups satisfies
    the axioms for a cohomology theory, namely:
    \begin{enumerate}[(i)]
        \ii If $f\simeq g\colon X\to Y$, then
	$f^* = g^*\colon H^n(Y)\to H^n(X)$.
	\ii For each CW pair $(X,A)$ there is a long exact sequence
	\begin{equation*}
	    \cdots
	    \overset\delta\too H^n(X/A)
	    \overset{q^*}\too H^n(X)
	    \overset{i^*}\too H^n(A)
	    \overset\delta\too H^{n+1}(X/A)
	    \overset{q^*}\too\cdots.
	\end{equation*}
	Equivalently, $H^n(X)\cong H^{n+1}(\Sigma X)$
	and $H^n(X/A)\to H^n(X)\to H^n(A)$ is exact for all $X$ and $n$.
	\ii Each $H^n$ sends coproducts to products,
	i.e. for a wedge $X = \bigvee_\alpha X_\alpha$
	with inclusions $i_\alpha\colon X_\alpha\to X$, the map
	\begin{equation*}
	    \prod_\alpha i_\alpha^*\colon H^n(X)\to\prod_\alpha H^n(X_\alpha)
	\end{equation*}
	is an isomorphism.
    \end{enumerate}
\end{theorem}

\begin{remark}
    The set $[X,K_n]$ has an (abelian) group structure
    for each $X$ due to $K_n \simeq \Omega^2 K_{n+2}$,
    with the group operation $f+g$ being
    $(f+g)(x) = f(x)\cdot g(x)$,
    the composition of the two maps $S^2\to K_{n+2}$
    given by $f(x)$ and $g(x)$.
\end{remark}

\begin{proof}
    (i) If $f$ and $g$ are homotopic,
    then for any homotopy class $[\gamma]\in [Y,K_n]$,
    \begin{equation*}
	f^*[\gamma] = [\gamma\circ f] = [\gamma\circ g] = g^*[\gamma].
    \end{equation*}

    (ii) TODO\todo{uses the cofibration sequence
    $A\to X\to X/A\to\Sigma A\to\cdots$ or something}

    (iii) The map into the product is invertible
    by sending a collection of maps $f_\alpha\colon X_\alpha\to K_n$
    to the map which is equal to $f_\alpha$ on
    each copy of $X_\alpha$ inside $X$.
\end{proof}

\section{Cohomology theories are represented by $\Omega$-spectra}

The converse of \Cref{thm:omega-spectra-cohomology} is also true ---
every cohomology theory is represented by an $\Omega$-spectrum!
To prove this, we first use a representability result
for a single nice enough functor:
\begin{theorem}[Brown representability theorem]
    \label{thm:brown-representability}
    Let $H\colon\mathbf{CW}^*\to\mathbf{Set}^*$ be a
    contravariant functor from
    the category of connected pointed CW complexes to
    the category of sets with the following properties:
    \begin{enumerate}[(i)]
        \ii (Homotopy invariance) If $f\simeq g\colon X\to Y$ then
	$H(f)\simeq H(g)\colon H(Y)\to H(X)$.
	\ii (Mayer-Vietoris axiom) Suppose $X$ is
	the union of (connected) subcomplexes $A$ and $B$.
	If there are $a\in H(A)$ and $b\in H(B)$ such that
	$a|_{A\cap B} = b|_{A\cap B}$ then
	there exists $x\in H(X)$ such that $x|_A = a$ and $x|_B = b$
	(where the restriction is the pullback under inclusion).
	\ii (Wedge axiom) For a wedge $X = \bigvee_\alpha X_\alpha$
	with inclusions $i_\alpha\colon X_\alpha\to X$, the map
	\begin{equation*}
	    \prod_\alpha i_\alpha^*\colon H(X)\to\prod_\alpha H(X_\alpha)
	\end{equation*}
	is an isomorphism.
    \end{enumerate}
    Then, there exists $K\in\mathbf{CW}^*$ and $u\in H(K)$
    such that $T_X\colon[X,K]\to H(X)$ given by $T_X(f) = f^*u$
    is an isomorphism for all $X$, i.e.
    $[-,K]$ and $H$ are naturally isomorphic.
\end{theorem}

\todo{why need connected?}

\begin{remark}
    The proper setting for this theorem should be
    the homotopy category of connected pointed CW complexes or something,
    in which case the functor is representable
    in the sense of being naturally isomorphic
    to a contravariant Hom functor (in $\mathbf{Set}^*$).
\end{remark}

The idea of the proof is that we construct a pair $(K,u)$
that satisfies the theorem statement when $X$ is a sphere
and then show that this turns out to also work for all spaces.

\begin{definition}[$n$-universal, $\pi_*$-universal]
    \label{def:n-pi-universal}
    A pair $(K,u)$ with $K\in\mathbf{CW}^*$ and $u\in H(K)$
    is \vocab{$n$-universal} if the transformation
    $T^{u}_{S^i}\colon\pi_i(K)=[S^i,K]\to H(S^i)$
    defined in the statement of \Cref{thm:brown-representability}
    is an isomorphism for $i < n$ and surjective for $i = n$.
    The pair is \vocab{$\pi_*$-universal} if it is $n$-universal for all $n$,
    i.e. $T_{S^i}$ is an isomorphism for all $i$.
\end{definition}

\begin{remark}
    \label{rmk:universal-htpy-eq}
    The word ``universal'' is because of the following:
    If $(K,u)$ and $(K',u')$ are $\pi_*$-universal
    and $f\colon K\to K'$ is a map such that $f^*u' = u$,
    then $T^{u'}_{S^n} \circ f_* = T^u_{S^n}$
    where $f_*\colon \pi_n(K)\to \pi_n(K')$
    is the induced homomorphism on homotopy groups,
    so $f$ is a bijection, hence an isomorpism.
    This implies that $f$ is a weak homotopy equivalence,
    which we can upgrade to a homotopy equivalence
    because $K$ and $K'$ are CW complexes.
\end{remark}

\begin{lemma}
    \label{lem:pi-universal-construction}
    For any pair $(Z,z)$ with $Z\in\mathbf{CW}^*$ and $z\in H(Z)$,
    there exists a $\pi_*$-universal pair $(K,u)$
    having $Z$ as a subcomplex and $u|_Z = z$.
\end{lemma}

\begin{proof}
    Let $K_1 = Z\vee\bigvee_\alpha S^1_\alpha$ over all $\alpha\in H(S^1)$.
    By the wedge axiom, there exists $u_1\in H(K_1)$ such that
    $u_1|_Z = u$ and $u_1|_{S^1_\alpha} = \alpha$,
    so $(K_1,u_1)$ is $1$-universal.

    The idea in general, given $n$-universal $(K_n,u_n)$,
    is to take maps $f_\alpha\colon S^n\to K_n$,
    one representing each element of the kernel of $T^{u_n}_{S^n}$,
    and form the mapping cone for $\bigvee_\alpha f_\alpha$
    to kill the kernel elements, and then to
    wedge with one $(n+1)$-sphere for each element of
    $H(S^{n+1})$ to get surjectivity.
    Then, to pass to the limiti (to get $n$-universality for all $n$),
    form the mapping telescope of the $K_n$s.
    \todo{details}
\end{proof}

\begin{lemma}
    \label{lem:pi-universal-extension-property}
    Let $(K,u)$ be a $\pi_*$-universal pair
    and let $(X,A)$ be a pointed $CW$ pair.
    Suppose $x\in H(X)$ and that $f\colon A\to K$ with $f^*u = x|_A$.
    Then $f$ extends to a map $g\colon X\to K$ with $g^*u = x$.
\end{lemma}

\begin{proof}
    Construct the reduced mapping cylinder of $f$,
    \begin{equation*}
	M_f = K\cup (A\times[0,1])/(f(a)\sim (a,0), (x_0,t)\sim(k_0));
    \end{equation*}
    where $x_0$ and $k_0$ are the basepoints of $X$ and $K$;
    since $M_f\simeq K$ we can regard $u$ as an element of $H(M_f)$
    (more precisely, there is a $\tilde u\in H(M_f)$
    which restricts to $u$) whose restriction to $A\times\{1\}$
    is equal to $x|_A$, since the inclusion sending $A$ to $A\times\{1\}$
    is homotopic to $f$. This means we can work with the inclusion
    $a\mapsto (a,1)$ of $A$ into $M_f$ instead of with $f$.

    Then, attach $X$ to $M_f$ along $A\times\{1\}$ to form $Z$.
    By the Mayer-Vietoris axiom there exists $z\in H(Z)$
    which restricts to $x$ and to $u$,
    and by \Cref{lem:pi-universal-construction}
    there is a $\pi_*$-universal pair $(K',u')$
    with $Z$ as a subcomplex of $K'$ such that $u'$ restricts to $z$.
    The argument of \Cref{rmk:universal-htpy-eq} now tells us that
    $K'\simeq M_f$, so $M_f$ deformation retracts onto $M_f$,
    inducing a homotopy from the inclusion of $X$ into $K'$
    to a map $g\colon X\to M_f$. This homotopy is constant on $A$,
    so $g$ extends the inclusion of $A$ into $M_f$,
    and $g^*u = g^* j^* u' = u'|_X = x$
    where $j\colon M_f\into K'$,
    by homotopy invariance of $H$.
\end{proof}

\begin{proof}
    [Proof of \Cref{thm:brown-representability}]
    \label{pf:brown-representability}
    We will show that a $\pi_*$-universal pair $(K,u)$ does the job.
    Taking $A$ to be a point in \Cref{lem:pi-universal-extension-property},
    note that $H(A)$ consists of a single point,
    say by applying the wedge axiom to $B\vee A$
    for any space $B$ with $H(B)$ nontrivial;
    this means $f^*u = x|_A$ holds for all $x\in H(X)$,
    so for each $x\in H(X)$ there is some $g_x\colon X\to K$
    with $g_x^*u = x$.

    Now, suppose that $T^u_X[f_0] = T^u_X[f_1]$
    for some $f_0,f_1\colon X\to K$ and let $p\colon X\times I\to X$
    be the projection (where $X\times I$ is the reduced cylinder,
    with $\{x_0\}\times I$ collapsed to a point).
    Define $x$ to be $p^*f_0^*u = p^*f_1^*u$;
    since projection is a left inverse to the inclusions
    of $X$ into the cylinder, $x|_{X\times\{0\}}$ and $x|_{X\times\{1\}}$
    are both equal to $f_0^*u = f_1^*u$.
    Therefore $x|_{(X\times\{0\})\vee(X\times\{1\})}$
    is equal to $(f_0^*u, f_1^*u)$. By
    \Cref{lem:pi-universal-extension-property},
    $f = f_0\vee f_1\colon (X\times\{0\})\vee (X\times\{1\})\to K$
    extends to $X\times I$, so $f_0$ and $f_1$ are homotopic,
    proving that $T^u_X$ is injective.
\end{proof}

\end{document}

